\begin{musica}{Amélia Dos Olhos Doces}{Carlos Mendes}{Tuna Universitária do Instituto Superior Técnico}
\nota{Introdução: [D A6 D A6 Em7 A7 Em7 A7] (2x) D A6}

A\ac{D}{mé}lia \ac{A6}{}__ dos olhos \ac{D}{do}ces, \ac{A6}{}__
\ac{D}{Quem} foi que te trouxe \ac{Ebdim7}{grá}vida de espe\ac{Em7}{ran}ça? \ac{A7 Em7 A7}{}__ __ __
Um \ac{Em7}{gos}to \ac{A7}{}__ de flor na \ac{Em7}{bo}ca, \ac{A7}{}__
Na pele e na \ac{Em7}{rou}pa per\ac{A7}{fu}mes de \ac{D}{Fran}ça. \ac{A6}{}__ _\ac{D}{}_ \ac{A6}{}__

\nota{Refrão:}
Ca\ac{D}{be}los \ac{A6}{cor} de vi\ac{D}{ú}va, \ac{A6}{}__
Cabelos de \ac{D7}{chu}va, sapatos de \ac{G}{ti}ras e pões,
Quantas \ac{Gm}{ve}zes não queres e não \ac{D}{a}mas \ac{Ebdim7}{}__
Os homens que \ac{Em7}{dor}mem,
Os homens que \ac{A7}{dor}mem contigo na \ac{D}{ca}ma. \ac{A6}{}__ _\ac{D}{}_ \ac{A6}{}__
Amélia dos olhos doces,
Quem dera que fosses apenas mulher;
Amélia dos olhos doces,
Se ao menos tivesses direito a viver.

\nota{Refrão}

A\ac{F#}{mé}lia gai\ac{Bm}{vo}ta, amante, po\ac{E}{e}ta, rosa de ca\ac{A}{fé};
A\ac{F#}{mé}lia ga\ac{Bm}{ia}ta, do bairro da \ac{E}{la}ta, do Cais do So\ac{A}{dré}.
Tens um nome de navio,
Teu corpo é um rio onde a sede corre
Olhos doces, quem diria,
Que o amor nascia onde Amélia morre

\nota{Refrão}

A\ac{D}{mé}lia \ac{A6}{}__ dos olhos \ac{D}{do}ces, \ac{A6}{}__
A\ac{D}{mé}lia \ac{A6}{}__ dos olhos \ac{D}{do}ces, \ac{A6}{}__
A\ac{Gm}{mé}lia \ac{A7}{dos} meus olhos \ac{D}{do}ces. \ac{Dmaj7}{}__

\posicoes{D: xx0232, D7: xx0212, Dmaj7: xx0222, Ebdim7: xx1212, E: 022100,
  Em7: 022030, F#: 244322, G: 320033, Gm: 355333, A: x02220, A6: x0x222,
  A7: x02020, Bm: x24432}
\end{musica}
